\section{Introduction}

Representation learning routinely stores internal states as vectors or tensors.
That storage choice does not specify which entities persist, which relations
compose, which metric quantities matter, how ordered and mass observables behave,
or how an action changes the state. TSI begins with a narrower question: what
mathematical data must be exposed if an internal state is intended to support
structural simulation, reasoning, and action?

The formal object developed here is a finite presented schema together with a
common tagged carrier, typed relations, a simplicial complex, a metric, a
filtration, a full-support mass, and action-conditioned tracked transitions.
These layers are allowed to interact, but they are not conflated. A vector may
encode the data; its coordinates are not the complete structural specification.

The paper makes one integrated contribution through five components:
(1) a state contract with a declared cross-layer compatibility axiom;
(2) an exact structural isomorphism and its equivalence proof;
(3) survivor-level preservation for five observables and its composition proof;
(4) a six-defect zero criterion separating local preservation from turnover; and
(5) categorical and quotient-level extensions, including metric--measure
comparison and rollout stability.

The finite-relation category, presented-category descent, filtered persistence
stability, metric--measure gluing, and the rollout recurrence are imported
mathematical ingredients. The TSI-specific contribution is their integration
under a common tagged state contract, the compatibility invariant, the
label-preserving partial transition category, and the explicit separation of
survivor preservation from temporal identity. The paper does not claim that
these layers are necessary or sufficient for intelligence, that they describe
human mathematical imagery, or that a neural model will discover them without
inductive constraints.
